\section{Jadro}

Jadro je hlavná časť práce a jeho členenie je určené typom práce. Vo vedeckých a odborných prácach má jadro spravidla tieto hlavné časti:

\begin{itemize}
	\item súčasný stav riešenej problematiky doma a v zahraničí,
	\item cieľ práce,
	\item metodika práce a metódy skúmania,
	\item výsledky práce,
	\item diskusia,
\end{itemize}

\subsection{Súčasný stav riešenej problematiky doma a v zahraničí}

V časti Súčasný stav riešenej problematiky autor uvádza dostupné informácie a poznatky týkajúce sa danej témy. Zdrojom pre spracovanie sú aktuálne publikované práce domácich a zahraničných autorov.

Podiel tejto časti práce má tvoriť približne 30 \% práce.

Záver kapitoly Súčasný stav riešenej problematiky podrobne popisuje problém, ktorý nebol doteraz v danej téme výskumu riešený, alebo bol riešený iným prístupom a nami predpokladané riešenie by mohlo priniesť lepšie výsledky.

\subsection{Ciele práce}

Časť Cieľ práce jasne, výstižne a presne charakterizuje predmet riešenia. Súčasťou sú aj rozpracované čiastkové ciele, ktoré podmieňujú dosiahnutie cieľa hlavného.

Príklad spracovania cieľa

Z analýzy súčasného stavu metodologických prístupov k vývoju IS a na základe analýzy prístupov k riešeniu transformácie CIM do PIM v MDA možno považovať za základný problém riešenia posunutie vývoja IS na čo najvyššiu úroveň abstrakcie. Dôvodom, prečo vývoj IS posúvať na čo najvyššiu možnú úroveň abstrakcie, je potreba čo najlepšej špecifikácie používateľských potrieb pre návrh informačného systému. Použitie princípov MDA prístupu pri návrhu IS slúži na vytvorenie základného piliera na preklenutie priepasti medzi doménou biznis analytikov a používateľov a doménou systémových analytikov a návrhárov. Zároveň v princípoch MDA architektúry je zárodok toho, ako povzniesť vývoj IS na čo najvyššiu úroveň abstrakcie, čo by mohlo viesť k odstráneniu závažných chýb, ktoré vznikajú pri tvorbe IS v súčasnosti.

Preto sme si stanovili cieľ práce: 

\textbf{Vytvoriť systémový prístup k transformácii CIM do PIM v MDA, podľa ktorého bude realizovaná transformácia modelu biznis procesov do vybraných návrhových modelov IS.}

Dosiahnutie tohto cieľa vyžaduje rozčlenenie riešenia do štyroch parciálnych cieľov: 

\begin{enumerate}
	\item Existencia viacerých notácií na reprezentáciu formálneho zápisu modelu biznis procesov vnáša do riešenia problém výberu najvhodnejšej notácie. Preto prvý parciálny cieľ práce je:
	
	\textbf{Výber notácií umožňujúcich formálny zápis modelov biznis procesov tak, aby bola zrozumiteľná pre všetky zainteresované strany}
	
	\item Úroveň PIM je reprezentovaná návrhovými modelmi v UML. Štandard UML poskytuje množstvo rôznych modelov na rôzne účely. Pre použitie v CIM – PIM transformácii nie sú všetky vhodné. Druhý parciálny cieľ je:
	
	\textbf{Špecifikácia UML modelov reprezentujúcich úroveň PIM.}
	
	\item Automatizovaná alebo polo – automatizovaná transformácia medzi modelmi nižších úrovní MDA je realizovaná na základe ich formálneho zápisu. Automatizovaná/polo – automatizovaná transformácia z CIM do PIM bude možná vtedy, keď budeme vedieť špecifikovať potrebný formálny zápis modelu biznis procesov. Pod formálnym zápisom je uvažovaný štandardizovaný grafický i sémantický (textový) formalizmus. To je úloha tretieho parciálneho cieľa:
	
	\textbf{Určenie požadovaného štandardizovaného formalizmu pre automatizovaný typ CIM – PIM transformácie.}
	
	\item Nájdenie vhodného problému a modelovanie jednotlivých modelov je podklad pre praktické overenie navrhovaných transformačných postupov. Overenie však bude iba čiastočné, lebo práca sa zameriava iba na najvyššie úrovne abstrakcie a transformácie medzi nimi. Pre úplné overenie by bolo potrebné previesť transformácie až po samotný zdrojový kód.
	
	\textbf{Posledným parciálnym cieľom je vytvorenie prípadovej štúdie, ktorá bude overením nami navrhovaného prístupu k CIM – PIM transformácii v MDA}
\end{enumerate}

\subsection{Metodika práce a metódy skúmania}

\begin{itemize}
	\item charakteristiku objektu skúmania, 
	\item pracovné postupy, 
	\item spôsob získavania údajov a ich zdroje,
	\item použité metódy vyhodnotenia a interpretácie výsledkov,
	\item štatistické metódy.
\end{itemize}

\textbf{Poznámky pre spracovanie kapitoly}

Metodika je podľa Wikipedie –Náuka o spôsobe vedeckej práce. Je to súhrn pracovných spôsobov a metód v určitej oblasti. Je to uvedomelý a cieľavedomý postup, určitým spôsobom usporiadaná činnosť alebo usporiadanie operácií, pretvárajúcich východiskové danosti istej cieľavedomej činnosti na jej zamýšľaný (čiastočne alebo úplne realizovaný) cieľ. 

Záverečné práce sú súčasťou výskumnej a vývojovej činnosti univerzity, ktorou sa realizuje vedecká práca. 

Spracovanie metodiky je závislé od:

\begin{itemize}
	\item riešenej témy,
	\item mechanizmu riešenia,
	\item použitých metód,
	\item spôsobu overenia,
\end{itemize}

Vzorom môže byť základný životný cyklus vývoja systémov s fázami: analýza/rozbor, návrh a implementácia/realizácia.

\textbf{Dôležitosť metodiky je v jej užitočnosti. Pre riešiteľa má byť hlavne:}

\begin{itemize}
	\item Spôsobom, ako systémovo a zmysluplne postupovať v riešení.
	\item Vzbudiť dôveru k dosiahnutiu cieľov.
	\item Odstrániť pocit  ťažkosti/prísnosti vedeckých prístupov. 
\end{itemize}

Ak riešiteľ aj napriek vyššie uvedeným dôvodom nepotrebuje alebo nevie vytvoriť metodiku vedeckovýskumnej práce, nemusí ju spracovať ešte pred riešením.  Metodika riešenia je však dôležitá pri záverečnom zhodnotení výsledkov. Treba ju dodatočne spracovať, aj keď postup riešenia bol ad hoc. Je potrebná pre posudzovanie správnosti riešenia je aj východiskom pre riešenie podobných výskumných problémov.

Metóda je podľa Wikipedie Metóda je poznaný zákon premenený na pravidlo, súbor pravidiel, systém regulatívnych princípov.

Spracovanie použitých metód v dizertačnej práci môže byť založené na nasledujúcej štruktúre:

\begin{itemize}
	\item Metódy vedeckého poznania, kde patria logické metódy
	\begin{itemize}
		\item analýza,
		\item syntéza,
		\item abstrakcia,
		\item konkretizácia,
		\item indukcia,
		\item dedukcia,
		\item \dots
	\end{itemize}
	
	\item Metódy vedeckého skúmania
	\begin{itemize}
		\item empirické skúmanie (pozorovanie, meranie, experiment, \dots),
		\item logické spracovanie empirických údajov (klasifikácia, typologizácia, korelácia, \dots),
		\item \dots
	\end{itemize}
\end{itemize}

\subsection{Výsledky práce a diskusia}

Výsledky práce a diskusia sú najvýznamnejšími časťami záverečnej práce. Výsledky (vlastné postoje alebo vlastné riešenie vecných problémov), ku ktorým autor dospel, sa musia logicky usporiadať a pri popisovaní sa musia dostatočne zhodnotiť. Zároveň sa komentujú všetky skutočnosti a poznatky v konfrontácii s výsledkami iných autorov. Ak je to vhodné, výsledky práce a diskusia môžu tvoriť aj jednu samostatnú časť a spoločne tvoria spravidla 30 až 40 \% záverečnej práce.

Podľa uváženia spracujte zvlášť kapitolu Výsledky práce a zvlášť Diskusia výsledkov, kde bude podrobný popis postupov podľa metodiky riešenia.

Výsledky, ktoré boli dosiahnuté riešením.

Diskusia výsledkov hlavne vtedy, ak budete porovnávať výsledky práce s inými riešeniami.

\subsection{Rozsah práce}

Odporúčaný rozsah bakalárskej práce je 30 až 40 strán (54 000 až 72 000 znakov vrátane medzier), diplomovej práce 50 až 70 strán (90 000 až 126 000 znakov), dizertačnej práce 80 až 120 strán (144 000 až 216 000 znakov) a habilitačnej práce do 150 strán. Strany sa počítajú od úvodu po záver. 


Podľa normy STN ISO 690 „Do rozsahu práce sa počítajú nasledovné časti textu: úvod, hlavný text, záver, zoznam bibliografických odkazov, citácie, poznámky pod čiarou